%% Section 03 — Background and Literature Review
\chapter{Background and Literature Review}
\label{ch:litreview}

\section{Introduction}
\label{sec:lr-intro}
Consensus is a collective decision-making process where every party can object to any
proposal and effectively veto it. International organisations are the standing forums through which states
take collective decisions. One becomes a consensus regime when consensus is
the only means to reach them. Importantly, consensus is not the same as unanimity: a unanimity
rule puts a proposal to a vote and records who votes in favour, against or abstain, whereas under
consensus a decision proceeds whenever no party registers a formal objection \parencite{Rietig2023}. 
Since a single party can halt any proposal, an outside observer might expect consensus regimes to
decide very little. In practice, \textcite{Sykorabodie2019} show that consensus regimes can reach agreement even on contested,
transnational issues, relying on informal practices (trust, dialogue, leadership) that fill the
gaps left by their formal rules. But success is never guaranteed and even historically productive regimes can stall \parencite{Gardiner2025}. 
Quantitatively, the picture is no different: what separates a consensus regime that keeps reaching decisions from one that can
no longer reach them is not clearly explained by existing models of consensus dynamics \parencite{Perez2018,Tran2024}. A model that better reproduces
how consensus is reached or broken is therefore needed, both to explain past behaviour and to anticipate when a
regime approaches the limits of its capacity to act.

Models of multi-actor decision-making already exist, but they were not built to explain why a consensus regime stops reaching decisions: they
represent each actor with a single, low-dimensional vector \parencite{DeGroot1974,Tsebelis2002} and rarely use the actual text those
actors produce \parencite{Botelho_CoalitionDynamics,Bernardo2021}, which is where they state what they want. To reproduce how a consensus regime behaves, a model needs positions across many issues,
read directly from the text the actors produce, and validated against observed outputs rather
than assumed. Recent advances in natural language processing and machine learning
\parencite{Devlin2019,Grootendorst_BERTopic} make each of these upgrades feasible.

The Antarctic Treaty system (ATs) is a consensus-driven regime. At its centre is the
Antarctic Treaty (AT), whose consultative parties meet yearly at the Antarctic Treaty
Consultative Meetings (ATCM), where outputs are adopted only by consensus \parencite{ATS2023RoP}.
Historically, the ATCM has been productive in terms of decisions reached. However, in 2022, for
the first time, the parties failed to reach consensus on the final report, and even as activity
has increased, the regime has become less responsive to challenges \parencite{Gardiner2025}. The
ATs is therefore a consensus regime approaching the limits of its capacity to act, and it is the
case study of this thesis.


\section{Modelling consensus: responsive and anticipatory approaches}
\label{sec:consensus}
There are two traditional groups of models for how groups reach consensus. The first group is responsive and interactive: each actor 
follows an update rule with no objective at all: it repeatedly reacts to
the current positions of the others (each position being that actor's opinion expressed as a number). These are the opinion-dynamics models.
 They treat consensus as \textit{convergence}: opinions are iteratively averaged until
they meet \parencite{DeGroot1974}, with several variants (\hyperref[app:e]{Appendix~E}).
Opinion-dynamics models can nevertheless test for agreement: declare consensus when all final
positions lie within a tolerance radius, as \textcite{Bernardo2021} do for the negotiation of the
2015 Paris Agreement. That test, however, captures numerical convergence. It does not model a proposal
that is entirely unacceptable to a party. In the ATs, any party can block a decision no matter how close
the remaining positions are.

The second group of models is anticipatory: each actor reasons about how the others will respond and
chooses its actions to maximise an objective function. These models use game theory to treat consensus as a
unanimity rule: in the veto-players framework of \textcite{Tsebelis2002}, the status quo is
replaced only when no party is left worse off (formal condition in \hyperref[app:e]{Appendix~E}). Each party accepts
only the policies it prefers to the status quo, so the set of adoptable policies is the
intersection of these per-party sets. As parties multiply and their positions spread, the sets
overlap less and the intersection shrinks \parencite{Tsebelis2002}.

Each group of models has strengths and limits. Responsive models are usually low-dimensional, can be scaled to many agents,
and are easy to calibrate against data. But their agents have no objective and no
strategy. Anticipatory models represent vetoes and strategic choice directly, and they explain \emph{why} a proposal passes or
fails. But their equilibria become hard to compute as parties and issues multiply, and they are
rarely confronted with data from real negotiations.

The ATs is this thesis's case study of such a regime: its Consultative Parties (CPs) take all
decisions by consensus at the ATCM.

Between these two groups lies a middle ground: models that retain the strategic, veto-bearing
structure of the second group but learn equilibrium from repeated interaction instead of solving
for it. This field is treated in the next section.


\section{Negotiation and coalition formation simulation using reinforcement learning}
\label{sec:marl}

The limits set out in Section~\ref{sec:consensus} imply two requirements for a model of ATs behaviour. 
First, the model must represent many heterogeneous parties, each holding a veto and bargaining over a status quo (i.e.
if no proposal passes, the current rules stay in force, so each party judges a proposal against them). Second, it
must also remain computable with that many parties and issues. Those two requirements cannot be tackled with the closed-form
 equilibria of analytic game theory. Multi-agent reinforcement learning (MARL) meets both requirements: it keeps the
game-theoretic substance but learns equilibrium behaviour from repeated interaction instead of
solving for it.

Formally, MARL is defined over stochastic (Markov) games, a generalisation of repeated games in
which the state of the negotiation changes as the agents act. Each agent observes the state,
chooses an action, and receives a reward: the objective function it learns to maximise. The
learning procedures approximate equilibrium behaviour over many simulated rounds rather than
deriving it in closed form \parencite{Littman1994}.

The MARL literature spans foundational bargaining and coalition-formation work
\parencite{Gronauer2022,Sarkar2022} and, most recently, \textcite{Kulkarni2025} uses agents based on Large Language Models
(LLMs) to model political negotiation directly (see \hyperref[app:e]{Appendix~E} for more applications). However, MARL simulations run on
synthetic data or abstract benchmarks\footnote{Standard test environments such as the board game
Diplomacy.} \parencite{Bolleddu2025} rather than the deliberative text records of a real consensus
regime.


\section{Quantitative and text-based study of institutions and treaties}
\label{sec:computational}

The quantitative study of international institutions long predates its application to
Antarctica. One line of work codes the design of treaties and organisations, such as their
decision rules, as structured variables, and links these variables to institutional performance,
survival, and decline \parencite{Gray2018}. A more recent trend treats the \emph{text} of
international agreements and debates as data \parencite{Grimmer2013}, recovering state positions
from voting records and speeches \parencite{Baturo2017,Slapin2008}. Both suggest that the
preferences and design features of treaty regimes can be measured, not merely described.

The ATs's health has been assessed qualitatively \parencite{Mancilla2023} and quantitatively \parencite{Gardiner2025,Botelho_CoalitionDynamics,Kumar_ConsistentPatterns}, in
two ways. \textcite{Gardiner2025} measures the regime in aggregate, counting inputs, outputs, and
the time legally binding Measures take to enter into force, and concludes that responsiveness to
challenges in the Antarctic region is declining. The second way is structural. A Working Paper
(WP), the document type through which parties bring proposals to the ATCM (\hyperref[app:f]{Appendix~F}), may be
submitted by a single party or co-authored by several, and this co-authorship has been used to
map which parties collaborate with which through co-authorship networks of ATCM documents
\parencite{Botelho_CoalitionDynamics,Kumar_ConsistentPatterns}.
This thesis uses co-authorship methods as validation targets: parties with similar stances
(directional positions on a topic) are likely to co-author papers together, and the same measures should help predict the output series
reported by \textcite{Gardiner2025} (Section~\ref{sec:ch2}).
Neither analysis, however, reads the content of the papers, so neither can say which issues
divide the parties, a limitation taken up in Section~\ref{sec:gap}.

A separate field does read the content: positional measures place each actor on a policy space
using only the text the actor produces. It has two paradigms, which differ in how they turn text
into numbers: word counts (i.e. frequencies) or embeddings, vectors in which texts with similar
meaning end up close together. By contrast, the word-count paradigm recovers a single, uninterpreted
dimension, which in a corpus as topically diverse as the ATCM Working Papers would likely
reflect what parties talk about rather than what they support. This thesis adopts the
embedding paradigm: topics are discovered with a pipeline \parencite{Grootendorst_BERTopic}, and
stances are read with natural language inference (NLI) classifiers, which judge whether a text
supports or contradicts a stated hypothesis \parencite{Burnham2023,Abercrombie2022}. The result
is a topic-conditioned stance: a separate position per party and topic through time. The
supporting literature and definitions are reviewed in \hyperref[app:g]{Appendix~G}, and Section~\ref{sec:ch2} sets
out the validation checks.


\section{Summary of literature gap}
\label{sec:gap}

The literatures reviewed above differ in method, but none recovers a directional preference: which country
is for or against what, extracted from the regime's own words (their limitations are collected
in \hyperref[app:e]{Appendix~E}).

This gap has two parts. The knowledge gap is the question raised in Section~\ref{sec:lr-intro}:
whether the regime's outputs, including its failures, can be anticipated from CPs'
positions at all. The methodological gap is
the missing instrument: no existing method measures each party's directional stance, per topic
and over time, from a consensus regime's own textual records, the ATs included. Without that instrument the question cannot be
answered. Section~\ref{sec:rqs} turns this gap into three research goals.

The closest prior work shares parts of this design but not its object. Stance detection recovers
for-or-against positions from parliamentary debates \parencite{Abercrombie2022}, but not a
per-party, per-topic measure through time. \textcite{Moghimifar2024} drive a negotiation
simulation with preferences derived from party manifestos, but in a parliamentary coalition
setting, not a consensus regime where every party holds a veto.
No prior work extracts a directional preference for each party over time from a consensus
regime's own textual records, nor uses such preferences to drive a strategic simulation in which agents have veto
power. This thesis supplies both halves and joins them: the preferences recovered from the record become the
 parameters of the simulation that reproduces consensus and its failures.


\section{Research problem}
\label{sec:rqs}
The first goal of this thesis is to reconstruct parties' preferences, so the regime's decisions
and failures can be anticipated from them. These preferences are extracted from the written
records as a directional measure (stance) across topics and time. The second goal is to extract
 directional opposition (veto) signals from the same records, and to test whether the support and
opposition signals together recover known history and predict decisions. The third goal is to use these preferences
to parameterise a multi-agent model that simulates consensus and its failures.


